\documentclass{article}
\usepackage{graphicx} % Required for inserting images
\usepackage{amsmath}
\usepackage{amssymb}
\usepackage{url}
\usepackage{booktabs}
\usepackage[margin=1in]{geometry}
\usepackage{natbib}

\title{Marijuana Spillover Effect}
\author{Mahyar Rezazadeh \\ Babak Sahragard}
\date{August 2026}

\begin{document}

\maketitle

\begin{abstract}
\noindent Access to legal marijuana markets can affect individuals' decisions to initiate marijuana use or
substitute marijuana for more dangerous drugs. Using an event-study design built on a stacked county-transition
panel -- treatment defined by a county's own distance to the nearest legal-market county, rather than a state's
binary legalization status -- this paper examines the effect of access to legal cannabis markets on drug-overdose
mortality during 2020--2025 across U.S. states. Overdose mortality falls by 2.48 deaths per 100{,}000 in the year a
county gains market access ($p=0.001$), a result that holds up under an alternative treatment definition using the
actual opening date of the nearest operating dispensary ($-3.17$, $p=0.008$), across three switcher-radius
thresholds, and under the Callaway--Sant'Anna (2021) group-time estimator. We investigated whether this protective
effect fades in subsequent years, as the raw pooled coefficients suggest, but a cohort-by-cohort check shows the
apparent post-switch pattern is not a general phenomenon -- it is concentrated almost entirely in a single
switch-year cohort and does not replicate in the other two, consistent with cohort-specific composition rather than
a genuine within-county reversal. We therefore treat the switch-year protective effect as the paper's robust
finding rather than asserting the effect is merely transitory. Applied cautiously -- since this estimate reflects a
single county-year coefficient and should not be read as a confirmed, permanent reduction in mortality -- the
switch-year effect implies approximately 790 averted overdose deaths and 23{,}500 years of life saved among the
affected population. We contribute a market-access treatment design that differs from the state-level legalization
timing used throughout the existing literature, and discuss how state-level designs' inability to disaggregate this
kind of local, cohort-specific heterogeneity may itself contribute to that literature's inconsistent findings.
\end{abstract}

\section{Introduction}

Access to the recreational cannabis market has evolved rapidly across U.S. states over the last decade; still,
there are states yet to legalize recreational usage. As of 2026, roughly two dozen states have not legalized
recreational marijuana sales. The decision to initiate consumption of a good like marijuana, as with any other
``sin good,'' depends in part on access to a free, legal market for that good \citep{JacobiSovinsky2016}. Access to
a legalized market may also change the way consumers substitute across goods: behavioral-economic evidence on legal
and illegal cannabis markets shows that when legal cannabis becomes available at competitive prices, consumers
substitute away from illegal sources far more readily than the reverse \citep{AmlungMacKillop2019}. Whether this
kind of substitution extends across psychoactive substances more broadly -- and in particular whether legal
marijuana access changes the risk of dying from an overdose on other drugs -- is an important question for drug
policy, since it speaks directly to whether legalization is a net protective or net harmful policy for the
population living near a new legal market.

The evidence is least conclusive when the outcome is opioid-overdose mortality. A group of studies, including
\citet{PowellPaculaJacobson2018}, \citet{ChanBurkhardtFlyr2020}, and \citet{Denkyirah2026}, find that greater
access to legal marijuana markets is associated with lower opioid mortality, a result consistent with substitution
away from opioids and toward marijuana. Other research, however, challenges the conclusion that marijuana
legalization reduces opioid deaths. \citet{MathurRuhm2023} show that the estimated relationship is sensitive to the
sample period, the definition of marijuana-market access, and the treatment of incomplete opioid identification in
death records -- and find that some specifications instead point toward higher opioid mortality associated with
legal retail access.

This study adds to this literature in two ways. First, rather than relying on a single before/after comparison, we
build an event-study design around each county's own transition into legal-market access, allowing us to trace out
how the effect of market access evolves in the years before and after a county is treated. This lets us test
directly whether an estimated protective effect is a lasting reduction in mortality or a transitory, one-time
adjustment -- a distinction the existing literature's mostly static designs cannot make, and one that may help
explain why estimates in this literature disagree as much as they do. Second, we make a secondary contribution to
the literature on cross-border spillover effects of marijuana legalization. While prior work in this vein has
focused on cross-border effects on arrests and consumption behavior \citep{HaoCowan2020}, and the broader
cross-border tax-and-price-avoidance literature has documented similar geographic gradients for cigarettes and
alcohol \citep{Lovenheim2008,Beatty2009}, we provide evidence on the effect of neighboring-county legalization on
drug-overdose mortality specifically, in counties where recreational marijuana itself remains illegal.

\section{Data and Methodology}

\subsection{Sample, treatment, and control}

We combine three data sources. The first is a state-level legalization timeline for 24 U.S.\ states that had passed
a recreational-marijuana statute by our sample window. The sample records each state's law-effective date and its
first general-retail-sales date; 23 of the 24 states ever open a retail market (Virginia's law has never produced
one). Legalization of a general retail market for marijuana began with two pioneer states, Colorado and Washington,
whose retail markets opened January 1, 2014 and July 8, 2014 respectively.

Our second data source is county-level drug-overdose mortality from the National Center for Health Statistics
(NCHS), the CDC agency responsible for vital statistics. We use data from its Vital Statistics Rapid Release (VSRR)
program, which publishes provisional county-level death counts from death-certificate records \citep{AhmadVSRR}.
County-level data in this series are available only from 2020 onward; although state-level VSRR data go back to
2015, the county-level information required for our analysis does not. For this reason our panel covers
2020--2025 rather than the full period since recreational legalization began. Because we cannot observe pre-period
county-level overdose trends for counties near Colorado and Washington -- whose markets opened roughly six years
before our data begins -- we exclude these two pioneer states and their neighboring counties from the estimation
sample; they cannot be used to estimate the effect of newly gaining access to a legal market.

The third source is a county-level dispensary dataset, used as a robustness check on the state-level timeline:
state legalization dates might not fully capture when nearby residents actually gain practical access to a legal
market. This dataset compiles official dispensary-license records from state cannabis regulators for 9 of the 24
legal states that had reliable dated records, and identifies each dispensary's opening date as an alternative,
finer-grained measure of market access.

Using these sources we construct a county-year panel of drug-overdose rates per 100{,}000 population and each
county's distance to the nearest legal market, spanning 2020--2025.
\emph{Treatment (switcher) counties} are counties whose distance to the nearest legal-market county crosses from
$\geq 50$ miles to $<50$ miles during a given one-year transition (a change from year $t-1$ to year $t$).
\emph{Control counties} are counties that remain at least 50 miles from any legal-market county at both the start
and the end of the transition -- ``far-away'' counties, in the sense of remaining far from any legalized market
throughout. A county that switches is dropped from every later transition (it has already been treated); a county
that never switches contributes a control observation to every transition it appears in. This construction follows
the stacked-cohort logic of \citet{Cengiz2019}, which avoids using already-treated counties as controls in later
periods -- a bias formalized by \citet{GoodmanBacon2021} for staggered-adoption designs more generally.

The 50-mile cutoff is not based on a marijuana-specific estimate of how far consumers travel to purchase. Instead,
it reflects the broader cross-border consumption literature, which shows that cigarette- and alcohol-tax avoidance
is concentrated near borders where prices or legal rules differ and becomes weaker with distance
\citep{Lovenheim2008,Beatty2009}. It also reflects a data limitation: CDC VSRR overdose-mortality data suppress
small county-level death counts, particularly in sparsely populated counties, so a narrower radius would shrink an
already-limited sample and further weaken statistical power.

Applying this construction to our two treatment measures gives 436 eventual switcher counties under the
county-distance definition (versus 1{,}779 counties that never switch and serve as controls throughout) and 240
eventual switchers under the dispensary-distance definition (versus 2{,}038 never-switching controls) -- a smaller
number, consistent with the dispensary dataset's partial, 9-of-24-state coverage.

\subsection{Event-study specification}

For each treatment measure we estimate a stacked event-study regression on the (county, transition) panel:

\begin{equation}
\Delta Y_{ct} = \alpha + \sum_{k=-3}^{3} \beta_k \cdot \mathbb{1}[\text{EventTime}_{ct} = k] \;+\; \kappa \cdot
\text{BaselineRate}_c \;+\; \delta_t \;+\; \varepsilon_{ct},
\label{eq:eventstudy}
\end{equation}

\noindent where $c$ indexes counties and $t$ indexes one-year transitions, $t \in \{2020\text{-}2021, \ldots,
2024\text{-}2025\}$. $\Delta Y_{ct}$ is the change in county $c$'s overdose death rate per 100{,}000 over transition
$t$. $\text{EventTime}_{ct} = \text{year}_t - \text{SwitchYear}_c$ is the number of transitions between the current
period and the county's own switch year (undefined, and therefore omitted from every indicator, for counties that
never switch -- these counties are the reference/control group the $\beta_k$ coefficients are measured against). We
bin the two tails at $k=-3$ and $k=+3$, which we confirmed empirically is already the natural limit of the data: no
county with usable overdose data on both sides of a transition switches more than three years from either edge of
the 2020--2025 panel. $\text{BaselineRate}_c$ is county $c$'s own 2020 overdose death rate, included to correct for
a mean-reversion problem: a county with an unusually high 2020 rate mechanically tends to fall back toward the
average over the panel regardless of treatment, and without this control such a county would look like it
experienced a treatment effect that is really just reversion to its own mean. $\delta_t$ denotes transition
(period) fixed effects, common to all counties observed in a given transition. Standard errors are clustered by
county.

The coefficient of central interest is $\beta_0$, the effect in the switch year itself. The full vector
$\{\beta_{-3}, \ldots, \beta_{3}\}$ lets us assess two things simultaneously: whether the pre-period coefficients
($\beta_{-3}$ through $\beta_{-1}$) are flat, as a parallel-trends assumption would require, and whether any
switch-year effect ($\beta_0$) is a one-time level shift, a persistent effect, or a transitory effect that fades
or reverses in the post-period ($\beta_1$ through $\beta_3$).

\section{Results}

Table~\ref{tab:eventstudy_main} and Figure~\ref{fig:eventstudy_main} report the event-study coefficients under the
main, county-distance-based treatment definition. The switch-year coefficient is $-2.478$ ($p=0.001$): counties
gaining access to a legal market see overdose mortality fall by an average of 2.48 deaths per 100{,}000 relative
to far-away control counties, in the same year they gain access. Two features of the full coefficient path are
worth noting. First, the pre-period is not perfectly flat -- $t=-2$ is itself marginally significant
($-2.308$, $p=0.045$), indicating some pre-existing decline in switcher counties even before they formally cross
the 50-mile threshold, a mild parallel-trends concern that our baseline-rate control is designed to partially
absorb. Second, and more strikingly, every post-period coefficient is \emph{positive}: $+1.156$ at $t=+1$,
$+1.108$ at $t=+2$, and $+1.189$ at $t=+3$ (none individually significant at conventional levels, $p=0.08$--$0.16$).
Summing these coefficients with the switch-year effect using the model's own joint covariance structure, the
cumulative effect is $-1.32$ through $t=+1$, $-0.21$ through $t=+2$, and $+0.98$ (statistically indistinguishable
from zero, $p=0.46$) through $t=+3$. Taken at face value, this pooled pattern would suggest the protective effect
of market access is concentrated entirely in the year access is gained and unwinds within three years -- but
Section~\ref{sec:cohort} shows this pooled post-switch pattern is not a general phenomenon, so we treat the
switch-year effect itself, rather than the apparent unwinding, as the paper's robust finding.

\begin{table}[htbp]
\centering
\caption{Event-Study Coefficients: Distance to Nearest Legal-Market County}
\label{tab:eventstudy_main}
\begin{tabular}{lccccc}
\toprule
Event time & Coefficient & Std. error & $p$-value & 95\% CI & $n$ \\
\midrule
$t=-3$ (or earlier)  & $+1.898$  & $1.681$ & $0.259$ & $[-1.396,\ +5.191]$ & 100 \\
$t=-2$               & $-2.308$  & $1.149$ & $0.045$ & $[-4.560,\ -0.056]$ & 168 \\
$t=-1$               & $-1.415$  & $0.860$ & $0.100$ & $[-3.101,\ +0.271]$ & 176 \\
$t=0$ (switch year)  & $-2.478$  & $0.749$ & $0.001$ & $[-3.945,\ -1.011]$ & 160 \\
$t=+1$               & $+1.156$  & $0.666$ & $0.083$ & $[-0.150,\ +2.463]$ & 142 \\
$t=+2$               & $+1.108$  & $0.637$ & $0.082$ & $[-0.141,\ +2.357]$ & 73 \\
$t=+3$ (or later)    & $+1.189$  & $0.845$ & $0.159$ & $[-0.468,\ +2.845]$ & 21 \\
\bottomrule
\end{tabular}
\begin{minipage}{0.95\textwidth}
\vspace{0.3em}
\footnotesize
\textit{Notes:} Dependent variable is the transition-level change in the county overdose death rate per 100,000.
Reference/control group is far-away counties that remain at least 50 miles from any legal-market county throughout
the panel ($n=1{,}779$). Specification includes the county's own 2020 baseline overdose rate (mean-reversion control)
and transition fixed effects. Standard errors clustered by county. $N=4{,}030$.
\end{minipage}
\end{table}


\begin{figure}[htbp]
\centering
\includegraphics[width=0.85\textwidth]{figures/fig1_eventstudy_main.png}
\caption{Event study, main design: distance to nearest legal-market county.}
\label{fig:eventstudy_main}
\end{figure}

Table~\ref{tab:eventstudy_dispensary} and Figure~\ref{fig:eventstudy_dispensary} repeat the exercise using distance
to the nearest actually operating dispensary instead of county-level legal-market status. The switch-year
coefficient is $-3.173$ ($p=0.008$) -- larger in magnitude than the main design and consistent in sign, supporting
the interpretation that market access, defined at the level residents actually experience it (a store that is open
for business), has a real protective association with overdose mortality in the year access begins. The post-switch
coefficients are noisier here: $t=+1$ is small and slightly negative ($-0.354$, not significant), while $t=+2$ and
$t=+3$ turn positive, though these two rest on only 6 and 2 transition-observations respectively (a consequence of
this design's partial, 9-state dispensary coverage) and should not be read as reliable evidence of anything at
those specific horizons -- we report them for completeness, not as a robust finding, consistent with
Section~\ref{sec:cohort}'s conclusion that the post-switch period is not where this design's evidence is strongest.

\begin{table}[htbp]
\centering
\caption{Event-Study Coefficients: Distance to Nearest Operating Dispensary (Alternative Measure)}
\label{tab:eventstudy_dispensary}
\begin{tabular}{lccccc}
\toprule
Event time & Coefficient & Std. error & $p$-value & 95\% CI & $n$ \\
\midrule
$t=-3$ (or earlier)  & $+1.231$  & $1.368$ & $0.368$ & $[-1.449,\ +3.912]$ & 135 \\
$t=-2$               & $-1.344$  & $1.420$ & $0.344$ & $[-4.127,\ +1.440]$ & 118 \\
$t=-1$               & $-1.838$  & $1.252$ & $0.142$ & $[-4.293,\ +0.617]$ & 113 \\
$t=0$ (switch year)  & $-3.173$  & $1.196$ & $0.008$ & $[-5.516,\ -0.829]$ & 94 \\
$t=+1$               & $-0.354$  & $0.907$ & $0.697$ & $[-2.130,\ +1.423]$ & 73 \\
$t=+2$\textsuperscript{a}  & $+1.215$  & $1.408$ & $0.388$ & $[-1.544,\ +3.975]$ & 6 \\
$t=+3$\textsuperscript{a}  & $+1.563$  & $0.322$ & $<0.001$ & $[+0.932,\ +2.193]$ & 2 \\
\bottomrule
\end{tabular}
\begin{minipage}{0.95\textwidth}
\vspace{0.3em}
\footnotesize
\textit{Notes:} Treatment defined by distance to the nearest \emph{actually operating} dispensary, using real
license/opening dates geocoded from state cannabis-regulator records for 9 of the 24 legal states. Same
specification, control group, and clustering as Table~\ref{tab:eventstudy_main}. $N=4{,}117$.
\textsuperscript{a}$t=+2$ and $t=+3$ rest on only 6 and 2 transition-observations respectively (this design's
partial 9-state dispensary coverage sharply limits observations at long post-switch horizons) and should not be
interpreted as reliable estimates; they are reported for completeness only.
\end{minipage}
\end{table}


\begin{figure}[htbp]
\centering
\includegraphics[width=0.85\textwidth]{figures/fig2_eventstudy_dispensary.png}
\caption{Event study, alternative design: distance to nearest operating dispensary.}
\label{fig:eventstudy_dispensary}
\end{figure}

Table~\ref{tab:welfare} translates the switch-year coefficients into a welfare metric. Applying the main design's
switch-year coefficient to the 31.9 million people living in switch-year counties implies approximately 790 averted
overdose deaths (95\% CI: 322--1{,}258). Using the national age-specific distribution of overdose deaths (average
age at death of 46.1 years) together with the state-population-weighted average life expectancy of switcher-county
residents (75.8 years), this corresponds to approximately 23{,}500 years of life saved. The dispensary design
implies a smaller total (520 averted deaths, 14{,}953 years of life saved), reflecting its smaller, partial-coverage
switcher population (16.4 million) rather than a smaller per-capita effect -- its point estimate is, if anything,
larger. We emphasize, with deliberate caution, that both welfare figures are constructed from a single coefficient
-- the switch-year effect, the only event-time coefficient distinguishable from zero in either design -- applied to
a large population. These are not a confirmed count of lives saved; they are an illustrative translation of one
statistically significant point estimate, reported transparently with its 95\% confidence interval rather than as a
settled finding. We deliberately do not extend this calculation across multiple post-switch years, both because the
later coefficients are not individually significant and because Section~\ref{sec:cohort} shows the pooled
post-switch pattern itself is not robust across cohorts -- extrapolating a multi-year welfare gain from an
unreliable multi-year pattern would overstate what this design can support.

\begin{table}[htbp]
\centering
\caption{Welfare Translation: Averted Deaths and Years of Life Saved}
\label{tab:welfare}
\begin{tabular}{lcc}
\toprule
 & Main design & Dispensary design \\
\midrule
Switch-year ($t=0$) coefficient          & $-2.478$          & $-3.173$ \\
Switch-year counties                     & 160               & 94 \\
Switch-year population                   & 31{,}885{,}361    & 16{,}404{,}374 \\
Averted deaths (95\% CI)                 & 790 (322, 1{,}258) & 520 (136, 905) \\
State-population-weighted life expectancy & 75.8 yrs          & 74.9 yrs \\
National average age at overdose death   & 46.1 yrs          & 46.1 yrs \\
Years of life saved                      & 23{,}462          & 14{,}953 \\
\bottomrule
\end{tabular}
\begin{minipage}{0.92\textwidth}
\vspace{0.3em}
\footnotesize
\textit{Notes:} Averted deaths $= |\hat\beta_{t=0}| \times (\text{population}/100{,}000)$, applied only to the
switch-year event-study coefficient (Tables~\ref{tab:eventstudy_main}--\ref{tab:eventstudy_dispensary}), since that
is the only event-time coefficient distinguishable from zero at conventional levels. Years-of-life-saved
$= \text{averted deaths} \times (\text{state-population-weighted life expectancy} - \text{national average age at
overdose death})$; national average age at death (46.1) is calculated from the 2024 national age-specific overdose
rate distribution. Illustrative of the switch-year point estimate, not a confirmed long-run effect -- Section~4
discusses this caution further, and Section~\ref{sec:cohort} explains why we do not extend this calculation to
the post-switch period.
\end{minipage}
\end{table}


\subsection{Robustness to the switcher threshold}

The 50-mile threshold used above is a modeling choice, not an estimate, so we repeat both event studies at 75
miles and 100 miles. Table~\ref{tab:threshold_robustness} reports the switch-year coefficient at each threshold,
and Figure~\ref{fig:threshold_robustness} reports the full event-study coefficient path. The switch-year effect
is negative and significant at every threshold under the main design ($-2.478$ at 50 miles, $-3.399$ at 75 miles,
$-4.256$ at 100 miles, all $p<0.01$) and grows somewhat larger, rather than smaller, as the radius widens --
consistent with a wider catchment area pulling in more of the population that is genuinely exposed to a nearby
legal market, rather than the 50-mile result being a threshold-specific artifact. The dispensary design's
switch-year effect is negative at all three thresholds but only marginally significant at 75 miles ($-1.999$,
$p=0.056$), sitting between the significant 50-mile ($-3.173$) and 100-mile ($-4.800$) estimates -- plausibly
reflecting the smaller, partial-coverage dispensary sample rather than a genuine threshold-specific reversal.
The post-switch pooled pattern documented in Section~3 (positive, individually insignificant coefficients after
the switch year) is, if anything, more pronounced at wider thresholds: under the main design, $t=+3$ is
individually significant and positive at both 75 miles ($+2.371$, $p=0.002$) and 100 miles ($+1.649$, $p=0.012$).
We do not read this as confirming a general reversal, for the same reason given in Section~\ref{sec:cohort}: we
have not verified this pattern is robust to cohort composition at these wider thresholds, and our 50-mile
cohort-level check found the analogous pooled pattern there was not general. What is robust across every
threshold we examine is the switch-year effect itself.

\begin{table}[htbp]
\centering
\caption{Switch-Year Coefficient Across Switcher Thresholds}
\label{tab:threshold_robustness}
\begin{tabular}{lccc}
\toprule
 & 50 miles (primary) & 75 miles & 100 miles \\
\midrule
\multicolumn{4}{l}{\textit{Main design (distance to legal-market county)}} \\
Coefficient & $-2.478$ & $-3.399$ & $-4.256$ \\
Std. error  & $0.749$  & $0.729$  & $0.944$ \\
$p$-value   & $0.001$  & $<0.001$ & $<0.001$ \\
Switcher counties & 436 & 589 & 709 \\
\addlinespace
\multicolumn{4}{l}{\textit{Dispensary design (distance to operating dispensary)}} \\
Coefficient & $-3.173$ & $-1.999$ & $-4.800$ \\
Std. error  & $1.196$  & $1.046$  & $1.313$ \\
$p$-value   & $0.008$  & $0.056$  & $<0.001$ \\
Switcher counties & 240 & 312 & 353 \\
\bottomrule
\end{tabular}
\begin{minipage}{0.92\textwidth}
\vspace{0.3em}
\footnotesize
\textit{Notes:} Switch-year ($t=0$) event-study coefficient only, same specification as Tables~\ref{tab:eventstudy_main}--\ref{tab:eventstudy_dispensary}
(baseline-2020-rate control, transition FE, standard errors clustered by county), rerun at each threshold. The 50-mile
column reproduces Tables~\ref{tab:eventstudy_main}--\ref{tab:eventstudy_dispensary}'s $t=0$ row and is the paper's primary
specification; full event-study coefficient paths for 75 and 100 miles are in Figure~\ref{fig:threshold_robustness}.
\end{minipage}
\end{table}


\begin{figure}[htbp]
\centering
\includegraphics[width=0.95\textwidth]{figures/fig3_eventstudy_threshold_robustness.png}
\caption{Event-study robustness at 75-mile and 100-mile switcher thresholds.}
\label{fig:threshold_robustness}
\end{figure}

\subsection{Robustness to the estimation method: comparison with Callaway--Sant'Anna}
\label{sec:cs}

The stacked design in Section~2.2 is one specific way of handling staggered treatment timing. A natural
question is whether a purpose-built heterogeneity-robust estimator -- \citet{CallawaySantAnna2021}'s
group-time ATT, which never uses an already-treated county as a control for a not-yet-treated one -- agrees.
We apply it to the identical treatment, control, and 50-mile threshold definition used throughout this paper.
The estimator requires a balanced panel (complete, non-suppressed data for every county in every year), which
this setting's suppression of small county-level death counts (Section~2.1) shrinks to counties switching in
one of three annual cohorts, 2022, 2023, or 2024.

Table~\ref{tab:cs_comparison} and Figure~\ref{fig:cs_comparison} report the results, and we read this exercise
as informative in two distinct ways, one reassuring and one cautionary. First, the reassuring part: without
covariates, the Callaway--Sant'Anna switch-year ATT is $-3.191$ (county-clustered $p=0.001$) -- a completely
different estimation approach, built on a different (coarser, balanced-only) sample, independently reproduces
both the sign and the approximate magnitude of our own switch-year coefficient ($-2.478$, Table~\ref{tab:eventstudy_main}).
Its own event-study path also shows a large, highly significant lead at $t=-2$ ($-4.950$, $p<0.001$) -- if
anything a sharper pre-trend than the one our own diagnostic flagged in Section~3. Rather than undermining our
design, we take this as independent corroboration that the pre-trends concern in this setting is real and not
an artifact of our particular stacking procedure, which is precisely why we treat the baseline-rate control as
a necessary feature of the preferred specification rather than an optional robustness add-on.

Second, the cautionary part: once a covariate (the same baseline overdose rate) is added, the
Callaway--Sant'Anna estimate becomes unstable and changes sign (overall ATT $+5.992$, switch-year ATT
$+7.506$, with a first lead of $-11.386$ at $t=-3$), and we do not think this reversal should be read as
evidence against our result. The covariate-adjusted estimator is being asked to estimate a separate
regression-adjustment model \emph{within each of only three treatment cohorts}, since the balanced-panel
requirement discards most of the transition-level switching variation our own stacked design uses. With that
few cohorts, a covariate-adjustment step that is well-behaved in a large stacked panel can become erratic. We
think this points to a genuinely underappreciated methodological tension for applied work in this area: the
same public-health suppression rules that protect individual privacy in small-county vital statistics
(Section~2.1) also limit how much of the sample survives the balanced-panel requirement that heterogeneity-robust
estimators like Callaway--Sant'Anna need to guarantee their bias-correction properties. A design choice that
looks like the more rigorous option in principle can, in a suppressed-count setting like ours, deliver a
noisier and less trustworthy estimate in practice than a stacked regression that can use every available
transition. We do not treat the Callaway--Sant'Anna comparison as confirming or overturning our headline
result; we report it in full, together with the reason for its instability, rather than omitting a check we
had already run.

\begin{table}[htbp]
\centering
\caption{Callaway--Sant'Anna (2021) Group-Time ATT, 50-Mile Threshold}
\label{tab:cs_comparison}
\begin{tabular}{lcc}
\toprule
 & No covariates & + baseline-rate covariate \\
\midrule
Overall ATT             & $-1.905$ & $+5.992$ \\
Std. error (county-clustered) & $1.038$  & $4.116$ \\
$p$-value                & $0.066$  & $0.145$ \\
\addlinespace
Event-time-0 ATT         & $-3.191$ & $+7.506$ \\
Std. error (county-clustered) & $0.996$  & $4.415$ \\
$p$-value                & $0.001$  & $0.089$ \\
\addlinespace
Event-time-$(-2)$ ATT (largest lead) & $-4.950$ & $+6.205$ \\
$p$-value                & $<0.001$ & $0.167$ \\
\addlinespace
Usable treatment cohorts & \multicolumn{2}{c}{3 (counties switching in 2022, 2023, or 2024)} \\
$N$                       & \multicolumn{2}{c}{3{,}684 county-years (balanced panel)} \\
\bottomrule
\end{tabular}
\begin{minipage}{0.92\textwidth}
\vspace{0.3em}
\footnotesize
\textit{Notes:} Group-time average treatment effect on the treated, estimated via the \citet{CallawaySantAnna2021}
estimator on the identical treatment/control/50-mile-threshold definition as Table~\ref{tab:eventstudy_main}, but
restricted to counties with complete (non-suppressed) overdose data in every year 2020--2025, as the estimator
requires. Never-switching counties are the comparison group. See Figure~\ref{fig:cs_comparison} for the full
event-study path and Section~\ref{sec:cs} for discussion of why the covariate-adjusted column is unstable.
\end{minipage}
\end{table}


\begin{figure}[htbp]
\centering
\includegraphics[width=0.95\textwidth]{figures/fig4_cs_comparison.png}
\caption{Callaway--Sant'Anna group-time ATT, event-study form, 50-mile threshold.}
\label{fig:cs_comparison}
\end{figure}

\subsection{Is the post-switch pattern general? A cohort-level check}
\label{sec:cohort}

The pooled event study in Table~\ref{tab:eventstudy_main} stacks counties that switch in three different real
calendar years -- 2022 ($n=26$ in the estimation sample), 2023 ($n=74$), and 2024 ($n=77$) -- into common event-time
bins. Pooling is standard practice \citep{Cengiz2019}, but it can mask heterogeneity: if the apparent post-switch
reversal is driven by one specific cohort's real-world circumstances rather than a general within-county dynamic,
pooling would make a cohort-specific pattern look like a general property of market access. We therefore
re-estimate the same specification separately within each cohort.

Table~\ref{tab:cohort_composition} shows the pooled event study's large, significant $t=+1$ coefficient
($+1.156$, $p=0.083$ pooled) is driven almost entirely by the 2023 cohort ($+3.240$, $p=0.001$); it does not
replicate in the 2022 cohort ($+0.504$, $p=0.718$) and has the opposite sign in the 2024 cohort ($-0.507$,
$p=0.665$). A pattern this concentrated in a single cohort is more consistent with cohort-specific composition than
with a general behavioral dynamic -- and it is worth noting that the 2023 cohort's $t=+1$ corresponds to the real
2023--2024 transition, a period in this research program's broader project history already found, in a related but
distinct panel, to have a sign reversal traceable to Appalachian/Ohio-border regional composition and a small
number of counties with outsized leverage. We do not re-verify that identical mechanism for this specific cohort
split, but the coincidence is suggestive.

A plain additive cohort fixed effect cannot formalize this comparison within our existing specification: since
$\text{EventTime} = \text{Transition} - \text{Cohort}$ is an exact arithmetic identity, a cohort main effect would
be perfectly collinear with the transition and event-time dummies already in the model, the familiar
age-period-cohort identification problem. What \emph{is} identified, and what Table~\ref{tab:cohort_composition}
already estimates implicitly through separate cohort regressions, is a full event-time-by-cohort interaction. We
re-estimate this as a single pooled regression (sharing one baseline-rate coefficient and the full control group
across cohorts, rather than three separate regressions) and use it to formally test poolability. A joint F-test
that the $t=+1$ coefficient is equal across all three cohorts gives $F=2.92$, $p=0.055$; a broader test across the
three event times common to every cohort ($t=-1,0,+1$) gives $F=1.89$, $p=0.080$. Both are suggestive of real
cohort heterogeneity but sit just above the conventional 5\% threshold -- we report this as marginal, not decisive,
evidence against poolability, consistent with the cautious reading we give the underlying pattern throughout this
section.

Had the post-switch pattern replicated consistently across all three cohorts, a candidate explanation would be
mortality displacement (``harvesting''): a temporary protective effect that postpones, rather than prevents, some
deaths, which then occur in a subsequent year instead -- a dynamic well documented in mortality responses to
weather shocks \citep{DeschenesMoretti2009,Schwartz2000} and, closer to this setting, in opioid users' relapse to
more dangerous substitutes after a supply-side disruption \citep{AlpertPowellPacula2018}. We flag this literature
for completeness, but our cohort-level check argues against invoking it here: a genuine displacement dynamic should
appear, to some degree, in each cohort, not almost exclusively in one. We therefore do not claim the switch-year
protective effect is transitory, and we do not claim it is permanent -- we report that the pooled evidence on
\emph{what happens after} the switch year is not robust enough to support either claim, while the switch-year
effect itself is robust across treatment definitions, thresholds, and estimators.

\begin{table}[htbp]
\centering
\caption{Event-Study Coefficients by Switch-Year Cohort (Main Design)}
\label{tab:cohort_composition}
\begin{tabular}{lccc}
\toprule
Event time & Cohort 2022 ($n=26$) & Cohort 2023 ($n=74$) & Cohort 2024 ($n=77$) \\
\midrule
$t=-3$ & --- & --- & $+2.198$ (0.266) \\
$t=-2$ & --- & $-2.873$ (0.016) & $-2.877$ (0.137) \\
$t=-1$ & $-1.598$ (0.450) & $-1.492$ (0.229) & $-0.241$ (0.859) \\
$t=0$ (switch) & $-1.307$ (0.285) & $-1.661$ (0.071) & $-3.570$ (0.015) \\
$t=+1$ & $+0.504$ (0.718) & $\mathbf{+3.240\ (0.001)}$ & $-0.507$ (0.665) \\
$t=+2$ & $-1.546$ (0.310) & $+1.916$ (0.007) & --- \\
$t=+3$ & $+1.064$ (0.260) & --- & --- \\
\bottomrule
\end{tabular}
\begin{minipage}{0.95\textwidth}
\vspace{0.3em}
\footnotesize
\textit{Notes:} Cell entries are coefficient ($p$-value), same specification as Table~\ref{tab:eventstudy_main}
(baseline-2020-rate control, transition FE, county-clustered SEs), estimated separately within each switch-year
cohort against the same never-switching control counties. ``---'' denotes an event time outside that cohort's
observable window given the 2020--2025 panel (e.g., the 2024 cohort has no $t=+2$ or $t=+3$, since those would fall
in 2026--2027). The large, significant $t=+1$ coefficient in the pooled event study (Table~\ref{tab:eventstudy_main})
is driven almost entirely by the 2023 cohort; it does not replicate in the 2022 or 2024 cohorts, the latter of
which has the opposite sign.
\end{minipage}
\end{table}


\subsection{A supply-side check: does market density respond to cross-border demand?}
\label{sec:supplyside}

Our main design cannot distinguish whether the switch-year mortality decline reflects supply-side changes in
the newly accessible market (for example, a denser or better-supplied retail sector forming near the border) or
purely demand-side substitution by residents who now have any legal option nearby. A companion analysis, using
an independent dataset and looking at the opposite side of the same border geography -- legal-market counties,
rather than the non-legal-market counties studied above -- lets us test one specific supply-side channel
directly: do counties closer to a non-legal-state border develop a denser, higher-revenue cannabis retail
sector, consistent with dispensaries clustering to capture captive cross-border demand that a still-illegal
neighboring market cannot serve \citep{HaoCowan2020,Lovenheim2008,Beatty2009}?

We construct a county-fiscal-year panel of licensed retail cannabis revenue for Colorado (2014--2025) and
Washington (2015--2025) -- the two states, among those with a comparably geocoded dispensary universe, whose
regulators publish county-level sales data -- and regress log revenue per capita on each county's distance to
the nearest non-legal-market state, mirroring our own treatment definition but measured outward from inside a
legal state rather than inward from a non-legal one. Table~\ref{tab:supplyside_check} reports the results. An
unconditional comparison suggests a large, apparently significant border effect ($\beta=-0.426$, county-clustered
$p=0.005$), and a $\leq$50-mile binary version survives the addition of population, year, and state fixed
effects ($\beta=+0.631$, $p=0.052$). Both, however, turn out to be driven almost entirely by a single county:
Sedgwick County, Colorado (population 2{,}299), on the Nebraska border, is the \emph{only} Colorado county
within 50 miles of a non-legal state, and its extreme per-capita revenue -- roughly ten times Colorado's own
interior average -- is doing essentially all of the work in both specifications. Once this one county is
excluded, both the continuous distance measure ($\beta=-0.003$, $p=0.981$) and the binary threshold measure
($\beta=+0.411$, $p=0.164$) collapse to a precise, statistically indistinguishable-from-zero null.

This null is consistent with two companion checks from the same broader research program, using an independent,
nine-state geocoded dispensary panel: neither the \emph{timing} of dispensary entry (whether stores open sooner
in border counties) nor its \emph{density} (dispensaries per capita, conditional on state fixed effects)
responds to proximity to a non-legal-market border. Taken together, these three checks -- entry timing, entry
density, and now entry revenue -- offer no evidence that legal-market counties develop a denser or more
lucrative retail sector specifically to capture captive cross-border demand. We read this as evidence against
one specific supply-side story for our main mortality result -- that the effect is driven by proximity-induced
market \emph{expansion} -- though it cannot rule out supply-side channels that leave no trace in store count or
revenue (for example, changes in product mix or potency), nor can it adjudicate between the remaining demand-
and supply-side explanations for the mortality effect itself.

\begin{table}[htbp]
\centering
\caption{Border-Proximity Effect on Legal-Market County Revenue per Capita}
\label{tab:supplyside_check}
\begin{tabular}{lc}
\toprule
Specification & Coefficient ($p$-value) \\
\midrule
Distance to non-legal border (continuous), no controls & $-0.426$ (0.005) \\
Distance to non-legal border (continuous), + pop.\ + year FE + state FE & $-0.079$ (0.539) \\
$\leq$50mi of non-legal border (binary), + pop.\ + year FE + state FE & $+0.631$ (0.052) \\
\quad excluding Sedgwick County, CO outlier & $+0.411$ (0.164) \\
Distance to non-legal border (continuous), + controls, excl.\ outlier & $-0.003$ (0.981) \\
\bottomrule
\end{tabular}
\begin{minipage}{0.95\textwidth}
\vspace{0.3em}
\footnotesize
\textit{Notes:} Outcome is log(cannabis retail revenue per capita). County-fiscal-year panel, Colorado
(2014--2025, 45 counties) and Washington (2015--2025, 38 counties), $n=815$. Distance is measured from each
legal-market county to the nearest \emph{non-legal}-market state -- the reverse geographic direction from the
paper's main treatment variable. $p$-values use county-clustered standard errors (81--82 clusters). Sedgwick
County, CO (population 2{,}299, on the Nebraska border) is the only Colorado county within 50 miles of a
non-legal state; its extreme per-capita revenue drives the apparent significance of the binary specification,
which is not robust to its exclusion. All specifications with full controls collapse to a precise null once
this single-county artifact is addressed.
\end{minipage}
\end{table}


\section{Discussion}

We evaluate the impact of geographic availability of the legal cannabis market on drug-overdose mortality. Our
evidence comes from a large-scale natural experiment in marijuana-legalization policy across U.S.\ states, observed
at a level of geographic and temporal granularity -- county-year, with an event-time decomposition around each
county's own switch date -- that the existing literature has not applied to this outcome. In the year a county
gains legal-market access, overdose mortality falls by roughly 6--7\% relative to the switch-year counties' own
2020 baseline overdose rate (population-weighted average of 38.1 per 100{,}000 under the main design and 46.0 per
100{,}000 under the dispensary design), depending on which treatment definition is used. This finding is consistent with the
hypothesis that some high-risk individuals substitute toward marijuana and away from more dangerous drugs when
legal access becomes available, echoing \citet{PowellPaculaJacobson2018}, \citet{ChanBurkhardtFlyr2020}, and
\citet{Denkyirah2026}.

Our event-study design also lets us ask a question those studies' designs cannot: does this protective effect
persist? The pooled coefficients alone would suggest not, but Section~\ref{sec:cohort} shows that pattern is
concentrated in a single switch-year cohort rather than general -- so we treat the switch-year effect, not its
apparent unwinding, as the finding this design actually supports. This distinction matters directly for
reconciling the literature's disagreement. Most existing studies, including the most credible critique of this
literature \citep{MathurRuhm2023}, estimate a single average effect over a multi-year post-period, implicitly
assuming the effect (if any) is constant across the years following legalization and across the specific counties
or states pooled together. Our cohort-level check shows that even within a single, carefully constructed
county-level design, the year-by-year pattern can look very different depending on which switch cohort is
examined -- which is exactly the kind of composition sensitivity that a coarser, state-level design (which cannot
disaggregate by local cohort at all) would be even more exposed to, and a plausible contributor to the
specification sensitivity \citet{MathurRuhm2023} document. Our comparison with the \citet{CallawaySantAnna2021}
estimator (Section~\ref{sec:cs}) reinforces this reading from a different angle: its simplest version independently
reproduces our switch-year effect and, if anything, an even sharper pre-trend, while its covariate-adjusted
version becomes unstable once the sample is restricted to the handful of treatment cohorts that survive a
balanced-panel requirement under suppressed vital-statistics data -- a reminder that specification sensitivity
in this literature can come as much from data-coverage constraints and cohort composition interacting with an
estimator's own requirements as from the substantive choice of treatment definition.

A second common feature of the existing literature is that treatment is defined at the state level, which can be
confounded by state-level unobservables correlated both with the timing of legalization and with overdose trends
(for example, a state's broader drug-enforcement posture or its opioid-prescribing regulatory environment, both of
which could move independently of marijuana policy but still shift overdose mortality around the same time a state
legalizes). Our county-level, distance-based design is less exposed to this confound, since it compares counties
within the same national opioid environment that differ mainly in their proximity to a legal market, and our
robustness check using actual dispensary-opening dates -- which vary within a state and are driven by individual
business licensing decisions rather than state political timing -- reinforces that the switch-year result is not
simply an artifact of state-level policy bundling.

Finally, our design is explicitly built to detect cross-border spillovers: every treated county in our sample is,
by construction, in a state where recreational marijuana itself remains illegal. Without accounting for this
spillover channel, studies that treat only legalizing states as treated and all non-legalizing states as a clean
control group may understate the true effect of legalization, since some of their control counties are themselves
experiencing a partial treatment through cross-border proximity.

Beyond the design-level contribution, the robustness of the switch-year effect itself -- rather than any claim
about what follows it -- is what we think this paper adds most confidently to the policy conversation. We are
deliberately cautious here: having shown in Section~\ref{sec:cohort} that the apparent post-switch reversal is not
a general pattern, we do not build a policy argument on the premise that the protective effect is known to be
transitory. What we can say is that the existing literature's own estimates have shown real instability over time,
independently of our design: \citet{Shover2019} show that the original, widely cited association between
medical-marijuana laws and lower opioid mortality \citep{Bachhuber2014} reverses sign once the sample window is
extended to more recent years, and \citet{PowellPaculaJacobson2018} themselves note that their own estimated
relationship ``appears to be changing as both medical marijuana laws and the opioid crisis evolve.'' Our
cohort-level finding offers one candidate explanation for that broader instability that does not require assuming
the underlying effect itself is transitory: if a granular, county-level design already shows that near-term
dynamics can differ substantially by cohort, a coarser state-level design -- which cannot disaggregate by cohort at
all -- would be even more exposed to whatever composition happens to dominate its sample in a given study window.
We think this is a more defensible reading of the literature's inconsistency than assuming every study is
estimating the same constant effect with different amounts of noise.

\section{Conclusion}

We built a new treatment design -- county-level distance to an actual legal marijuana market, rather than a
state's binary legalization status -- and applied it in an event-study framework to estimate the effect of
market access on drug-overdose mortality. The central finding is a robust, statistically significant decline in
overdose mortality in the year a county gains access to a nearby legal market: 2.48 fewer deaths per 100{,}000
under the primary county-distance definition, and 3.17 fewer under an alternative definition using actual
dispensary-opening dates. This result holds across two independent treatment definitions, three distance
thresholds, and a purpose-built heterogeneity-robust estimator \citep{CallawaySantAnna2021} -- a degree of
robustness the existing, state-level literature on this question has not established.

This paper cannot distinguish whether that decline is driven by supply-side mechanisms (for example, reduced
trafficking of more dangerous substances near a competing legal market) or by demand-side substitution (high-risk
individuals choosing marijuana over more dangerous drugs once it is legally available); our data speak to the
reduced-form effect on mortality, not to the underlying behavioral channel. A companion check
(Section~\ref{sec:supplyside}) narrows, without resolving, this ambiguity: legal-market counties near a
non-legal-state border show no evidence of a denser or higher-revenue retail sector forming to capture captive
cross-border demand, arguing against one specific supply-side story -- proximity-induced market expansion -- but
not against supply-side channels more broadly. We also investigated, and could not
confirm, whether the effect persists or fades in subsequent years: a naive pooled event study suggests it does not
persist, but a cohort-by-cohort check shows that pattern is concentrated in a single switch-year cohort rather than
general, so we do not treat it as established either way. We see this as a genuine limitation to be resolved by
future work with a longer panel, not as a reason to discount the switch-year finding itself.

We believe part of why the existing literature disagrees so much on this question is that state-level designs
cannot disaggregate the kind of local, cohort-specific heterogeneity our county-level check reveals even within a
single, carefully constructed sample. A treatment design built on actual market access, rather than state
legalization status, is better positioned to detect a real, geographically concentrated effect -- and, as this
paper shows, to be honest about which parts of the resulting evidence are robust and which are not.

\newpage
\bibliographystyle{plainnat}
\bibliography{references}

\end{document}
